% Лекция 17. Источник: Тервер(Лекция)_202611112317_42212.pdf, с. 6.
\clearpage
\begingroup
\tcbset{myboxstyle/.append style={breakable=false}}
\everymath{\textstyle}
\subsection{Коэффициент корреляции}

\begin{definition}
    Пусть $0<\DD\xi<\infty$ и $0<\DD\eta<\infty$.
    \textbf{Коэффициентом корреляции} случайных величин $\xi$ и $\eta$
    называется число
    \[
        \operatorname{corr}(\xi,\eta)
        =\frac{\operatorname{cov}(\xi,\eta)}
        {\sqrt{\DD\xi\cdot\DD\eta}}.
    \]
\end{definition}

\begin{theorem}[Свойства коэффициента корреляции]
    При конечных положительных дисперсиях справедливы следующие свойства.
    \begin{enumerate}
        \item $|\operatorname{corr}(\xi,\eta)|\leq1$.
        \item Если $\xi$ и $\eta$ независимы, то
        $\operatorname{corr}(\xi,\eta)=0$.
        \item $\operatorname{corr}(\xi,\eta)=1$ тогда и только тогда,
        когда существуют $a>0$ и $b\in\R$, для которых
        \[
            \xi=a\eta+b\quad\text{почти наверное}.
        \]
        \item $\operatorname{corr}(\xi,\eta)=-1$ тогда и только тогда,
        когда существуют $a<0$ и $b\in\R$, для которых
        \[
            \xi=a\eta+b\quad\text{почти наверное}.
        \]
    \end{enumerate}
\end{theorem}
\endgroup
